% !TeX program = pdflatex
\documentclass[11pt,a4paper]{article}

% Compile with pdfLaTeX:
%   latexmk -pdf legged_mani_analysis_draft.tex
\usepackage[utf8]{inputenc}
\usepackage[T1]{fontenc}

\usepackage{amsmath,amssymb,amsthm}
\usepackage{bm}
\usepackage{geometry}
\usepackage{booktabs}
\usepackage{array}
\usepackage{algorithm}
\usepackage{algorithmic}
\usepackage{xcolor}
\usepackage{hyperref}

\geometry{margin=2.4cm}
\emergencystretch=1em
\hypersetup{
  colorlinks=true,
  linkcolor=blue!55!black,
  urlcolor=blue!55!black,
  citecolor=blue!55!black
}
\setcounter{tocdepth}{2}

\newcommand{\R}{\mathbb{R}}
\newcommand{\norm}[1]{\left\lVert #1 \right\rVert}
\newcommand{\sqnorm}[1]{\left\lVert #1 \right\rVert^2}
\newcommand{\clip}{\operatorname{clip}}
\newcommand{\diag}{\operatorname{diag}}
\newcommand{\sat}{\operatorname{sat}}
\newcommand{\code}[1]{\texttt{#1}}

\title{\code{legged\_mani}: B2--Z1 Push-Box Controller\\
Mathematical Formulation and Implementation Analysis}
\author{Notes based on analysis of the \code{legged\_mani} codebase}
\date{\today}

\begin{document}
\maketitle
\tableofcontents
\newpage

%=============================================================================
\section{Scope and Executive Summary}
%=============================================================================

This document analyzes the \code{push\_box} task implemented by
\path{control/controllers/mppi_push_box.py}. The controller moves a free box
toward a fixed world-frame goal with the end effector (EE) of a four-degree-
of-freedom Unitree Z1 arm mounted on a Unitree B2 quadruped. The analysis is
based primarily on:
\begin{itemize}
  \item push-box task logic and rollout cost:
        \path{control/controllers/mppi_push_box.py},
  \item shared sampling and whole-body support:
        \path{control/controllers/base_controller.py} and
        \path{control/controllers/whole_body_arm_controller.py},
  \item kinematic base-pose planning:
        \path{control/planning/base_pose_cem.py},
  \item task and controller parameters:
        \path{utils/tasks.py} and \path{configs/mppi_push_box.yml},
  \item rollout and simulation models:
        \path{models/b2_z1_base_push_box.xml} and
        \path{models/scene_push_box.xml}.
\end{itemize}

The controller combines two optimizers with different roles:
\begin{enumerate}
  \item A background cross-entropy method (CEM) searches a three-dimensional
        kinematic base-pose reference
        $[\Delta x_b,\Delta y_b,z_b]^\top$ that makes the moving box-contact
        point reachable.
  \item Whole-body MPPI samples 16-dimensional desired joint-position
        trajectories and evaluates their full robot--box contact dynamics in
        parallel MuJoCo rollouts.
\end{enumerate}
The CEM result is only a reference. MPPI remains the sole dynamic trajectory
optimizer. Its cost combines robot state tracking, virtual PD effort, box
position and tilt, moving EE-to-box contact tracking, and arm--torso
clearance. The contact point changes from a pre-contact point outside the box
to a small virtual penetration inside the pushing face after engagement.

%=============================================================================
\section{Robot, Box, State, and Input}
%=============================================================================

\subsection{Degrees of Freedom}

The push-box model contains the 23 robot position coordinates and 22 robot
velocity coordinates of the B2--Z1 system, plus a seven-coordinate free-joint
pose and six velocities for the box:
\begin{equation}
  n_q=30,\qquad n_v=28,\qquad n_u=16.
\end{equation}
The actuated coordinates are the twelve B2 leg joints and the first four Z1
arm joints:
\begin{align*}
\bm{q}_j = [&q_{\mathrm{FR,h}},q_{\mathrm{FR,t}},q_{\mathrm{FR,c}},
q_{\mathrm{FL,h}},q_{\mathrm{FL,t}},q_{\mathrm{FL,c}},\\
&q_{\mathrm{RR,h}},q_{\mathrm{RR,t}},q_{\mathrm{RR,c}},
q_{\mathrm{RL,h}},q_{\mathrm{RL,t}},q_{\mathrm{RL,c}},
q_{a1},q_{a2},q_{a3},q_{a4}]^\top .
\end{align*}
Here $\mathrm{h}$, $\mathrm{t}$, and $\mathrm{c}$ denote hip, thigh, and calf.

\subsection{Observation State}

The controller receives the complete MuJoCo state
$[\bm{q}_{\mathrm{pos}}^\top,\bm{q}_{\mathrm{vel}}^\top]^\top$ and separates
it into a canonical 45-dimensional robot state and a 13-dimensional box
state:
\begin{equation}
  \boxed{
  \bm{x}=
  \begin{bmatrix}
    \bm{x}_r\\
    \bm{x}_o
  \end{bmatrix}
  \in\R^{58},}
\end{equation}
\begin{equation}
  \bm{x}_r=
  \begin{bmatrix}
    \bm{p}_b\\
    \bm{q}_b\\
    \bm{q}_j\\
    \bm{v}_b\\
    \bm{\omega}_b\\
    \dot{\bm{q}}_j
  \end{bmatrix}
  \in\R^{45},
  \qquad
  \bm{x}_o=
  \begin{bmatrix}
    \bm{p}_o\\
    \bm{q}_o\\
    \bm{v}_o\\
    \bm{\omega}_o
  \end{bmatrix}
  \in\R^{13}.
\end{equation}
The base and box quaternions use MuJoCo's scalar-first convention
$[q_w,q_x,q_y,q_z]^\top$. The robot state is formed by excluding the box
free-joint indices from the full arrays; the box state is extracted using the
\code{box\_joint} position and velocity addresses. This is safer than assuming
that the two subsystems always occupy fixed contiguous blocks.

The MuJoCo FULLPHYSICS rollout buffer additionally stores time:
\begin{equation}
  [t,\bm{q}_{\mathrm{pos}}^\top,\bm{q}_{\mathrm{vel}}^\top]^\top
  \in\R^{59}.
\end{equation}
The leading time element is removed before cost evaluation, leaving the
58-dimensional observation above.

\subsection{Control Input and Actuator Model}

The optimized input is a vector of desired joint positions, not joint torques
or contact forces:
\begin{equation}
  \boxed{\bm{u}=\bm{q}_{j,\mathrm{des}}\in\R^{16}},
  \qquad
  u_i^{\mathrm{clip}}=\sat_{[u_i^-,u_i^+]}(u_i).
\end{equation}
Each affine MuJoCo position actuator then computes
\begin{align}
  \tau_i^{\mathrm{unsat}}
  &=k_{p,i}(u_i^{\mathrm{clip}}-q_{j,i})
    -k_{d,i}\dot q_{j,i},\\
  \tau_i^{\mathrm{MJ}}
  &=\sat_{[\tau_i^-,\tau_i^+]}
    \left(\tau_i^{\mathrm{unsat}}\right).
\end{align}
The hip actuators use $(k_p,k_d)=(100,5)$, the thigh and calf actuators use
$(300,8)$, and the four arm actuators use $(300,5)$. Desired positions are
clipped by the sampler, while MuJoCo applies the separate actuator-force
limits.

%=============================================================================
\section{Prediction Dynamics and Contact}
%=============================================================================

\subsection{Discrete MuJoCo Dynamics}

One prediction step is the full contact-rich transition
\begin{equation}
  \boxed{
  \bm{x}_{k+1}
  =F_{\mathrm{MJ}}(\bm{x}_k,\bm{u}_k;
                    \bm{\theta}_{\mathrm{model}}).
  }
\end{equation}
The transition includes floating-base and articulated rigid-body dynamics,
position-PD actuation, force saturation, ground contact, and physical
EE--box contact. The rollout model disables only physical contact between the
arm subtree and the torso collision geometry; their clearance is instead
handled by the analytic soft cost and hard rejection described later. Arm
contact with the box and ground remains enabled.

The configured integration interval and horizon are
\begin{equation}
  \Delta t=0.01\,\mathrm{s},\qquad H=40,\qquad
  T_H=H\Delta t=0.4\,\mathrm{s}.
\end{equation}
The rollouts use a pyramidal friction cone and the enabled contact override
\begin{equation}
  \code{o\_solref}=[0.02,1.0].
\end{equation}
The simulator copies the rollout model's cone, override flag,
\code{o\_solref}, \code{o\_solimp}, friction override, and margin so that
predicted and executed contact options match.

\subsection{Implicit Contact Scheduling}

No contact force appears in $\bm{u}$, and there is no prescribed stance,
swing, or EE--box contact mode in the dynamic optimization. Leg contact is
determined by the height-conditioned gait prior together with MuJoCo
collision. Box pushing emerges from physical arm--box contact when sampled
joint-position commands bring the EE into the box.

The Boolean \code{contact\_engaged} flag does not measure a MuJoCo contact
force. It switches when the measured EE position is sufficiently close to a
geometric pre-contact target. It changes the reference used by IK and the
cost, but physical contact is still resolved exclusively by MuJoCo.

%=============================================================================
\section{Finite-Horizon Stochastic Optimal Control Problem}
%=============================================================================

For a fixed reference supplied by the task logic and CEM planner, the intended
MPPI problem can be written as
\begin{equation}
\boxed{
\begin{aligned}
\min_{\bm{U}}\quad
  &J(\bm{U};\bm{x}_0)
   =\sum_{k=0}^{H-1}\ell_k(\bm{x}_{k+1},\bm{u}_k)
    +\Phi_{\mathrm{EE}}(\bm{x}_H),\\
\mathrm{s.t.}\quad
  &\bm{x}_{k+1}=F_{\mathrm{MJ}}(\bm{x}_k,\bm{u}_k),\\
  &\bm{u}_{\min}\leq\bm{u}_k\leq\bm{u}_{\max},\\
  &d_{k,s}\geq d_{\mathrm{hard}},\qquad s=1,2,3,\\
  &\bm{x}_0=\hat{\bm{x}}(t),
\end{aligned}}
\label{eq:finite_ocp}
\end{equation}
where $\bm{U}=[\bm{u}_0,\ldots,\bm{u}_{H-1}]$ and $d_{k,s}$ is the
arm-capsule-to-torso clearance. The box goal is a soft objective, not a hard
terminal constraint.

The controller samples $N=30$ candidate sequences around a structured
gait/IK prior. In parallel, the CEM planner periodically solves a separate
static kinematic problem for the base reference. CEM does not propagate
dynamics, score box motion over the MPPI horizon, or replace the MPPI
optimization in Eq.~\eqref{eq:finite_ocp}.

%=============================================================================
\section{Push Geometry, Base-Pose Planning, and Task Progression}
%=============================================================================

\subsection{Box Goal and Push Direction}

The single configured box-position target is
\begin{equation}
  \bm{p}_o^*=[3.0,2.0,0.19]^\top\;\mathrm{m}.
\end{equation}
For a predicted box position $\bm{p}_{o,k}$, the horizontal push direction is
\begin{equation}
  \bm{d}_k^{xy}
  =
  \frac{\bm{p}_o^{*,xy}-\bm{p}_{o,k}^{xy}}
       {\norm{\bm{p}_o^{*,xy}-\bm{p}_{o,k}^{xy}}},
  \qquad
  \bm{d}_k=
  \begin{bmatrix}\bm{d}_k^{xy}\\0\end{bmatrix}.
\end{equation}
When the horizontal distance is numerically zero, the most recent valid push
direction is reused.

\subsection{Moving Contact Target}

Let $\bm{R}_{o,k}$ be the world-from-box rotation and
$\bm{h}_o=[0.19,0.19,0.19]^\top\,\mathrm{m}$ the box half-size. The distance
from the box center to its surface along $\bm{d}_k$ is
\begin{equation}
  \rho_k
  =\min_{i:\,|(\bm{R}_{o,k}^\top\bm{d}_k)_i|>0}
  \frac{h_{o,i}}{|(\bm{R}_{o,k}^\top\bm{d}_k)_i|}.
\end{equation}
The moving EE target lies on the face opposite the desired motion:
\begin{equation}
  \tilde{\bm{p}}_{e,k}^*
  =\bm{p}_{o,k}-\rho_k\bm{d}_k+\delta\bm{d}_k,
\end{equation}
after which its vertical component is explicitly overwritten as
\begin{equation}
  p_{e,k}^{*,z}=p_{o,k}^{z}-0.05\,\mathrm{m}.
\end{equation}
The offset $\delta$ defines two reference modes:
\begin{equation}
  \delta=
  \begin{cases}
    -0.05\,\mathrm{m},&\text{before engagement},\\
    +0.015\,\mathrm{m},&\text{after engagement}.
  \end{cases}
\end{equation}
Thus the first target stays outside the rear face, whereas the second is a
small virtual penetration into the face that encourages sustained physical
contact. Every rollout recomputes this target from its own predicted box pose;
the cost does not track only the currently observed target.

Engagement occurs when the current EE is within $0.04\,\mathrm{m}$ of the
pre-contact target. The switch is one-way until success and triggers a CEM
refresh for the new penetration target.

\subsection{Success Condition}

Only horizontal box error determines task completion:
\begin{equation}
  \norm{\bm{p}_o^{xy}-\bm{p}_o^{*,xy}}
  \leq0.20\,\mathrm{m}.
\end{equation}
The inequality must hold for 20 consecutive calls to \code{goal\_reached};
the counter resets to zero whenever the box leaves the tolerance. On success,
pending CEM work is invalidated, contact engagement is cleared, the
\code{stance\_hold} gait is selected, and the controller holds the current
base position.

%=============================================================================
\subsection{Kinematic CEM Base-Pose Planner}
%=============================================================================

\subsubsection{Decision Variable and Sampling}

The asynchronous planner searches
\begin{equation}
  \bm{z}=
  [\Delta x_b,\Delta y_b,z_b]^\top
\end{equation}
using a private MuJoCo model, private IK data, and a dedicated random
generator. Its feasible search region is
\begin{equation}
  \norm{[\Delta x_b,\Delta y_b]^\top}\leq0.40\,\mathrm{m},
  \qquad
  0.35\leq z_b\leq0.543542\,\mathrm{m}.
\end{equation}
CEM draws 30 samples per iteration for three iterations, retains the best
20\% as elites, and exponentially smooths the new distribution with factor
$0.15$. The first iteration also includes deterministic seeds at three
forward offsets and three height levels. A complete plan normally evaluates
$30\times3=90$ kinematic candidates and launches no additional MPPI rollout.

\subsubsection{CEM Objective}

For a candidate base pose, position-only IK produces an arm posture and EE
residual $r$. Define
\begin{align}
  \eta_r&=\frac{r}{0.03},&
  v_r&=\max(\eta_r-1,0),\\
  \eta_{xy}&=\frac{\norm{[\Delta x_b,\Delta y_b]^\top}}{0.40},&
  \eta_h&=\frac{0.543542-z_b}{0.543542-0.35}.
\end{align}
Let $\ell_{\mathrm{clr}}$ denote the normalized arm--torso clearance deficit
and $\ell_{\mathrm{lim}}$ the normalized arm joint-limit margin deficit. The
implemented CEM score is
\begin{equation}
\begin{split}
  J_{\mathrm{CEM}}={}&
  \eta_r^2
  +100v_r^2
  +1000\,\mathbb{I}_{\mathrm{infeasible}}
  +6\eta_{xy}^2
  +3\eta_h^2\\
  &+2\ell_{\mathrm{clr}}
  +\ell_{\mathrm{lim}}
  +10^6\mathbb{I}_{\mathrm{arm\ collision}}.
\end{split}
\end{equation}
A CEM candidate is labeled feasible only when $r\leq0.03\,\mathrm{m}$ and
the arm--torso hard-clearance test passes.

Before a CEM pose is applied, the controller also checks planar base--box
separation. Using support radii of the oriented box and torso along their
center-to-center direction, it requires an additional $0.05\,\mathrm{m}$
clearance. A collision-free pose may be used as a
\code{fallback-intermediate} plan even if its IK residual exceeds the CEM
reachability threshold; such a plan is flagged for later replanning.

\subsubsection{Asynchronous Reference Management}

CEM runs on one background worker while MPPI continues using the current
reference. A completed result is discarded if the request ID is obsolete, if
the contact target moved more than $0.10\,\mathrm{m}$, or if the box moved
more than $0.10\,\mathrm{m}$ since submission. A refresh is requested when:
\begin{itemize}
  \item engagement changes the contact target,
  \item the active target drifts by at least $0.08\,\mathrm{m}$, or
  \item a fallback plan settles without becoming feasible.
\end{itemize}
Requests are separated by at least 10 controller steps, and fallback
replanning is limited to eight attempts.

The accepted base reference is approached with rate limits. Height changes at
at most $0.10\,\mathrm{m/s}$. Horizontal speed is at most
$1.0\,\mathrm{m/s}$ at standing height and decreases linearly to 40\% of that
value at the minimum height. Position and height tolerances use hysteresis to
avoid repeatedly switching between walking and settled modes.

%=============================================================================
\section{Nominal Whole-Body Trajectory}
%=============================================================================

\subsection{Height-Conditioned Leg Gait}

The default and tracking gait is \code{stance\_hold}; \code{walk\_fast} is
selected while planar base motion is required. Each gait has phase-compatible
references at base heights
\begin{equation}
  z_b\in\{0.35,\;0.45,\;0.543542\}\;\mathrm{m}.
\end{equation}
The scheduler linearly interpolates joint positions and velocities between
the two surrounding height levels. When the commanded height changes, the
joint-velocity reference also includes the morphing derivative induced by
$\dot z_b$. Gait switches preserve the previous phase instead of restarting
the gait.

\subsection{Arm IK Reference}

At every controller update, the active moving contact target is converted to
a four-joint arm reference with position-only damped least-squares IK:
\begin{align}
  \bm{e}^{(r)}
  &=\bm{p}_e^*-\bm{p}_e(\bm{q}_a^{(r)}),\\
  \Delta\bm{q}_a^{(r)}
  &=\bm{J}_a^\top
    \left(\bm{J}_a\bm{J}_a^\top+\lambda_{\mathrm{IK}}^2\bm{I}\right)^{-1}
    \bm{e}^{(r)},\\
  \bm{q}_a^{(r+1)}
  &=\clip\left(
    \bm{q}_a^{(r)}+\alpha_{\mathrm{IK}}\Delta\bm{q}_a^{(r)}
    \right).
\end{align}
The MPPI-side update uses damping $0.05$, step size $0.7$, and at most five
iterations. The new solution is filtered with smoothing factor $0.3$ and a
per-update joint-change limit of $0.04\,\mathrm{rad}$. The CEM reachability
test uses the same IK structure but allows 40 iterations.

\subsection{Whole-Body Prior and Residual Warm Start}

The horizon reference is
\begin{equation}
  \bm{r}_{j,k}=
  \begin{bmatrix}
    \bm{q}^{g}_{\ell,k}\\
    \bm{q}^{\mathrm{IK}}_a\\
    \dot{\bm{q}}^{g}_{\ell,k}\\
    \bm{0}_4
  \end{bmatrix}
  \in\R^{32}.
\end{equation}
Its position component plus a warm-started correction defines the nominal
action:
\begin{equation}
  \bar{\bm{u}}_k=
  \begin{bmatrix}
    \bm{q}^{g}_{\ell,k}\\
    \bm{q}^{\mathrm{IK}}_a
  \end{bmatrix}
  +\Delta\bm{u}^{\mathrm{warm}}_k.
\end{equation}
After optimization, the selected correction relative to the gait/IK nominal
is shifted by one step, and the new tail correction is set to zero. For the
first 50 controller updates, the gait/IK reference is blended from the
standing initialization.

%=============================================================================
\section{Sampling and MPPI Update}
%=============================================================================

\subsection{Rollout Sampling Modes}

The current controller exposes three main rollout modes:
\code{original}, \code{original\_spline}, and \code{safe\_spline}. Both the
controller constructor and \code{simulate\_mppi.py} select
\code{original\_spline} by default. All three modes use $K=4$ horizon-aligned
knots
\begin{equation}
  \mathcal{K}=\{\kappa_j\}_{j=0}^{K-1},
  \qquad
  \kappa_j=
  \operatorname{round}\left(\frac{j(H-1)}{K-1}\right),
\end{equation}
For $H=40$, this gives $\mathcal{K}=\{0,13,26,39\}$. The implementation
requires the rounded indices to be unique. Let
$\mathcal{S}_{\mathcal{K}}[\cdot]_k$ denote cubic-spline reconstruction at
horizon index $k$ from values defined at these knots. Each sample draws
\begin{equation}
  \bm{\epsilon}_{n,j}
  \sim\mathcal{N}(\bm{0},\diag(\bm{\sigma}^2)),
  \qquad j=0,\ldots,K-1.
\end{equation}

For the gait-aware modes, define the full-rate gait/IK nominal
$\bm{g}_k$, the shifted MPPI correction $\bm{c}_k$, and the clipped absolute
warm start
\begin{equation}
  \bm{a}_k
  =\clip(\bm{g}_k+\bm{c}_k,\bm{u}_{\min},\bm{u}_{\max}).
\end{equation}
The three sampling rules are as follows.

\paragraph{\code{original}.}
This mode disables the explicit gait/IK nominal
(\code{nominal\_from\_gait}$=$\code{False}) and maps internally to
\path{sample_type=cubic_original}. Here $\bm{a}_k$ denotes the
previous optimized absolute action sequence after receding-horizon shifting.
The knot actions are perturbed and the complete absolute action is
reconstructed:
\begin{equation}
  \boxed{
  \bm{u}_{n,k}^{\mathrm{original}}
  =\clip\left(
    \mathcal{S}_{\mathcal{K}}
    [\bm{a}_{\kappa_j}+\bm{\epsilon}_{n,j}]_k,
    \bm{u}_{\min},\bm{u}_{\max}
  \right).
  }
\end{equation}
After the MPPI update, the selected absolute sequence is shifted and its last
action is repeated at the new horizon tail.

\paragraph{\code{original\_spline}.}
This is the default mode and maps internally to
\path{sample_type=cubic_gait_absolute}. It enables the gait/IK nominal and
forms $\bm{a}_k$ from the full gait/IK/residual warm start, but then splines
that complete absolute trajectory:
\begin{equation}
  \boxed{
  \bm{u}_{n,k}^{\mathrm{original\_spline}}
  =\clip\left(
    \mathcal{S}_{\mathcal{K}}
    [\bm{a}_{\kappa_j}+\bm{\epsilon}_{n,j}]_k,
    \bm{u}_{\min},\bm{u}_{\max}
  \right).
  }
\end{equation}
Consequently, both the height-conditioned gait/IK component and the
warm-started correction are reconstructed from only four knots. Short
full-rate gait features between knots can therefore be smoothed.

\paragraph{\code{safe\_spline}.}
This mode maps internally to \path{sample_type=cubic_gait_residual}. It
preserves the complete full-rate gait/IK nominal and applies the cubic spline
only to the correction plus exploration noise:
\begin{equation}
  \boxed{
  \bm{u}_{n,k}^{\mathrm{safe\_spline}}
  =\clip\left(
    \bm{g}_k+
    \mathcal{S}_{\mathcal{K}}
    [\bm{c}_{\kappa_j}+\bm{\epsilon}_{n,j}]_k,
    \bm{u}_{\min},\bm{u}_{\max}
  \right).
  }
\end{equation}
It is therefore the only one of the three modes that preserves the nominal
gait/IK samples at every horizon step before final action clipping.

Candidate zero is overwritten with a mode-consistent noise-free baseline:
\begin{center}
\begin{tabular}{p{0.22\textwidth}p{0.62\textwidth}}
\toprule
Mode & Candidate-zero baseline\\
\midrule
\code{original}
& $\bm{u}_{0,k}=\clip(\bm{a}_k)$, the full shifted absolute warm start\\
\code{original\_spline}
& $\bm{u}_{0,k}=\clip(
   \mathcal{S}_{\mathcal{K}}[\bm{a}_{\kappa_j}]_k)$\\
\code{safe\_spline}
& $\bm{u}_{0,k}=\clip(
   \bm{g}_k+\mathcal{S}_{\mathcal{K}}[\bm{c}_{\kappa_j}]_k)$\\
\bottomrule
\end{tabular}
\end{center}
Thus every batch contains a zero-noise reference appropriate to the selected
sampling family. In \code{original\_spline} and \code{safe\_spline}, the
selected correction relative to $\bm{g}_k$ is shifted after optimization and
the new tail correction is set to zero.

The base exploration standard deviations are $0.03$ for each hip, $0.10$ for
each thigh and calf, and $0.06\,\mathrm{rad}$ for each arm joint. The
\code{stance\_hold} gait multiplies these values by $0.6$;
\code{walk\_fast} uses scale $1.0$.

\subsection{Cost Normalization and Weights}

Let $\mathcal{V}$ be the candidates that pass the hard arm--torso test and
have finite total cost. For $n\in\mathcal{V}$,
\begin{align}
  S_{\min}&=\min_{n\in\mathcal{V}}S_n,\\
  \Delta S&=\max_{n\in\mathcal{V}}S_n-S_{\min},\\
  w_n&=
  \begin{cases}
    \exp\!\left[
      -\dfrac{1}{\lambda_{\mathrm{MPPI}}}
       \dfrac{S_n-S_{\min}}{\Delta S}
    \right],
    &\Delta S>10^{-12},\\
    1,&\Delta S\leq10^{-12},
  \end{cases}
\end{align}
and $w_n=0$ for invalid candidates. The configured temperature is
$\lambda_{\mathrm{MPPI}}=0.1$. The updated sequence is the weighted average
of the absolute candidate actions:
\begin{equation}
  \boxed{
  \bm{u}^{+}_k=
  \clip\left(
    \frac{\sum_{n=1}^{N}w_n\bm{u}_{n,k}}
         {\sum_{n=1}^{N}w_n+10^{-10}},
    \bm{u}_{\min},\bm{u}_{\max}
  \right).
  }
\end{equation}
If every candidate is invalid, the last accepted safe trajectory is shifted
and its final action is repeated. Only $\bm{u}^{+}_0$ is executed.

%=============================================================================
\section{Cost Function}
%=============================================================================

\subsection{Robot State and Virtual PD Effort}

The robot reference is assembled from the ramped base pose and velocity
commands, the height-conditioned leg gait, and the arm IK reference. Before
cost evaluation, the actual world-frame base linear velocity is transformed
to the base-local frame.

Let $\bm{e}_{r,k}$ be the robot-state error after replacing the raw quaternion
difference with base roll and pitch errors and zeroing the remaining two
quaternion slots. The robot and virtual-effort cost is
\begin{align}
  \tilde{\bm{\tau}}_k
  &=\bm{K}_p(\bm{u}_k-\bm{q}_{j,k})
    -\bm{K}_d\dot{\bm{q}}_{j,k},\\
  \ell_{r,k}
  &=\bm{e}_{r,k}^{\top}\bm{Q}_r\bm{e}_{r,k}
    +\tilde{\bm{\tau}}_k^\top\bm{R}\tilde{\bm{\tau}}_k.
\end{align}
The nonzero diagonal weights are:
\begin{itemize}
  \item base $x/y/z$: $3000$, $3000$, and $1500$,
  \item base roll/pitch: $5000$ and $3000$,
  \item all 16 joint positions: $500$ each,
  \item base-local linear velocity $x/y$: $20$ each,
  \item base linear $z$ and all three angular velocities: $0.1$ each,
  \item virtual PD effort: $\bm{R}=10^{-3}\bm{I}_{16}$.
\end{itemize}
Yaw and all joint-velocity tracking weights are zero. The virtual effort does
not include MuJoCo actuator-force saturation and therefore need not equal the
force actually applied in the rollout.

\subsection{Box Position and Tilt}

The box position cost retains the weighted L1 form used by the push-box
controller:
\begin{equation}
  \boxed{
  \ell_{o,k}
  =1000|p_{o,k}^x-p_o^{*,x}|
   +1500|p_{o,k}^y-p_o^{*,y}|.
  }
\end{equation}
The configured vertical box-position weight is zero, and box linear and
angular velocities are not directly penalized.

Let $\bm{R}_{o,k}$ be the box rotation. Its yaw-invariant tilt angle is
\begin{equation}
  \theta_{o,k}=\arccos\!\left(
  \clip([\bm{R}_{o,k}]_{33},-1,1)\right).
\end{equation}
The orientation cost is
\begin{equation}
  \ell_{\mathrm{tilt},k}
  =10^4\theta_{o,k}^2
   +10^5\mathbb{I}\{\theta_{o,k}>0.35\}.
\end{equation}
The second term is a large soft discontinuous penalty; it does not mark a
rollout invalid.

\subsection{Moving EE Contact Cost}

For the rollout-dependent target derived in Section 5,
\begin{equation}
  \ell_{p,k}
  =7000\sqnorm{\bm{p}_{e,k}-\bm{p}_{e,k}^*}.
\end{equation}
The implementation also supports the sign-invariant quaternion error
\begin{equation}
  e_{R,k}
  =2\arccos\left(
  \clip(|\bm{q}_{e,k}^\top\bm{q}_e^*|,-1,1)\right),
  \qquad
  \ell_{R,k}=w_R e_{R,k}^2,
\end{equation}
but the push-box configuration sets $w_R=0$. EE orientation tracking is
therefore disabled.

The final EE term receives an additional multiplier of 10:
\begin{equation}
  \Phi_{\mathrm{EE}}(\bm{x}_H)
  =10(\ell_{p,H-1}+\ell_{R,H-1}).
\end{equation}
Because the same error is already present in the last stage, its total final
coefficient is 11. There is no analogous terminal multiplier on the box
position or tilt cost.

\subsection{Arm--Torso Clearance}

Four frames---shoulder, elbow, wrist, and EE---form three arm capsules with
radii
\begin{equation}
  [0.030,0.030,0.0375]\;\mathrm{m}.
\end{equation}
The torso is the oriented box associated with \code{base\_collision}, with
half-size $[0.25,0.14,0.075]\,\mathrm{m}$. The proximal $0.07\,\mathrm{m}$ of
the first capsule is excluded near the structural shoulder mount.

For signed capsule clearances $d_{k,s}$, the soft cost is
\begin{equation}
  \boxed{
  \ell_{c,k}
  =10^4\sum_{s=1}^{3}
  [\max(0.005-d_{k,s},0)]^2.
  }
\end{equation}
A complete rollout is valid only if
\begin{equation}
  d_{k,s}\geq0
  \qquad\text{for every }k\text{ and }s.
\end{equation}
Thus capsule penetration is rejected, while a tangent configuration is
allowed. The shared collision helper first uses uniformly sampled centerline
points to bound clearance and invokes the exact segment-to-box distance only
when needed for hard-validity classification.

\subsection{Total Stage and Rollout Cost}

The complete stage cost is
\begin{equation}
  \ell_k
  =\ell_{r,k}+\ell_{o,k}+\ell_{\mathrm{tilt},k}
   +\ell_{p,k}+\ell_{R,k}+\ell_{c,k}.
\end{equation}
Candidate $n$ has total cost
\begin{equation}
  \boxed{
  S_n
  =\sum_{k=0}^{H-1}\ell_k^{(n)}
   +10\left(
     \ell_{p,H-1}^{(n)}+\ell_{R,H-1}^{(n)}
   \right).
  }
\end{equation}
Only arm--torso hard clearance and non-finite cost affect the validity mask;
large box tilt remains part of $S_n$ rather than causing hard rejection.

%=============================================================================
\section{Complete Receding-Horizon Algorithm}
%=============================================================================

\begin{algorithm}[ht]
\caption{\code{push\_box} CEM-guided whole-body MPPI update}
\label{alg:push_box_mppi}
\begin{algorithmic}[1]
  \REQUIRE Observation $\hat{\bm{x}}$, active base plan, previous safe MPPI plan
  \STATE Split $\hat{\bm{x}}$ into robot and box states.
  \STATE Compute the box-to-goal direction and moving EE contact target.
  \STATE Engage the penetration target if the EE reaches the pre-contact target.
  \STATE Apply a completed non-stale CEM plan, if available.
  \STATE Submit a background CEM refresh when required and allowed by cooldown.
  \STATE Ramp the base reference toward the accepted CEM pose.
  \STATE Select \code{walk\_fast} during planar motion; otherwise select
         \code{stance\_hold}.
  \STATE Solve smoothed position IK for the current moving EE target.
  \STATE Combine height-conditioned gait and arm IK with the shifted MPPI
         correction.
  \STATE Generate $N=30$ candidates using the selected
         \code{original}, \code{original\_spline}, or
         \code{safe\_spline} rule.
  \STATE Preserve one mode-consistent noise-free baseline candidate.
  \STATE Roll out every candidate for $H=40$ MuJoCo steps with FK sensors.
  \STATE Evaluate robot, box, EE-contact, tilt, and arm-clearance costs.
  \STATE Set zero weight for every hard-collision or non-finite rollout.
  \IF{at least one valid candidate exists}
    \STATE Normalize valid costs and compute exponential weights.
    \STATE Take the weighted mean of absolute candidate actions and save it as
           the latest safe plan.
  \ELSE
    \STATE Shift and reuse the previous safe plan.
  \ENDIF
  \STATE Shift either the absolute plan (\code{original}) or the optimized
         gait/IK correction (gait-aware modes), and advance gait phase.
  \RETURN The first desired joint-position command $\bm{u}^{+}_0$
\end{algorithmic}
\end{algorithm}

The standard \code{simulate\_mppi.py} configuration uses a $0.01\,\mathrm{s}$
simulation step and a $100\,\mathrm{Hz}$ controller rate, so one MPPI update
and one task-success check are requested per simulation step.

%=============================================================================
\section{Default Configuration Summary}
%=============================================================================

\begin{center}
\begin{tabular}{>{\raggedright\arraybackslash}p{0.40\textwidth}
                >{\raggedright\arraybackslash}p{0.43\textwidth}}
\toprule
Parameter & Default value\\
\midrule
Robot / object & B2 with four-DoF Z1 / free $0.38\,\mathrm{m}$ box\\
Full state / input dimension & $58$ / $16$\\
Robot / box state dimension & $45$ / $13$\\
Input interpretation & desired joint position\\
Time step / horizon & $0.01\,\mathrm{s}$ / 40 steps ($0.4\,\mathrm{s}$)\\
MPPI samples / workers & 30 / 5\\
MPPI temperature & $0.1$\\
Rollout mode / knots & \code{original\_spline} (default) / four\\
Initial keyframe & \code{stand}\\
Tracking / approach gait & \code{stance\_hold} / \code{walk\_fast}\\
Height-conditioned levels & $0.35$, $0.45$, $0.543542\,\mathrm{m}$\\
CEM variable & $[\Delta x_b,\Delta y_b,z_b]^\top$\\
CEM samples / iterations / elite fraction & 30 / 3 / 0.20\\
CEM maximum horizontal shift & $0.40\,\mathrm{m}$\\
Box position target & $[3.0,2.0,0.19]^\top\,\mathrm{m}$\\
Pre-contact gap / push penetration & $0.05/0.015\,\mathrm{m}$\\
EE position/orientation weights & 7000 / 0\\
EE terminal scale & 10\\
Box $x/y/z$ L1 weights & 1000 / 1500 / 0\\
Box tilt weight / threshold / violation penalty & $10^4$ / $0.35$ rad / $10^5$\\
Collision safe/hard distances & $0.005/0.0\,\mathrm{m}$\\
Collision soft weight & $10^4$\\
Box success tolerance / hold & $0.20\,\mathrm{m}$ / 20 checks\\
\bottomrule
\end{tabular}
\end{center}

% The former Section 11, ``Comparison with OCS2 Legged MPC'', is intentionally
% omitted for the push-box analysis. Preserve the later section numbers so the
% omission is visible in the table of contents.
\setcounter{section}{11}

%=============================================================================
\section{Implementation Caveats and Limitations}
%=============================================================================

\begin{enumerate}
  \item \textbf{The update differs from canonical MPPI.}
  Costs are normalized by the valid batch's min--max range, and absolute
  candidate actions are averaged. There is no control-noise correction term.
  Consequently, $\lambda=0.1$ controls selection pressure over normalized
  relative cost rather than temperature in the original physical cost units.

  \item \textbf{Engagement is geometric, not contact-sensor based.}
  The target switches after an EE position-tolerance test. It does not verify
  normal force, sustained contact, or whether the EE is on the intended face.

  \item \textbf{The CEM planner is kinematic and local.}
  It searches at most $0.40\,\mathrm{m}$ from the base pose at request time,
  assumes an upright base, and does not predict balance, gait feasibility, box
  motion, or dynamic contact. A fallback pose may intentionally remain outside
  the configured IK residual threshold.

  \item \textbf{The moving contact target is a point objective.}
  EE orientation has zero weight. The point is placed at a fixed vertical
  offset from the box center and does not optimize contact normal alignment,
  friction cone margin, or wrench direction.

  \item \textbf{Box tilt is not a hard constraint.}
  Exceeding $0.35$ rad adds $10^5$ to the cost but does not invalidate the
  rollout. Whether this is sufficient depends on the scale and spread of all
  other accumulated costs.

  \item \textbf{The collision model is incomplete.}
  Hard rejection covers three arm capsules against one torso box. It does not
  guarantee arm--leg, gripper--environment, box--base, or general
  self-collision safety over MPPI rollouts. The separate CEM acceptance check
  covers only the planned base footprint against the current box footprint.

  \item \textbf{Analytic and physical collision models are intentionally
  different.}
  Arm--torso physical contact is disabled in the rollout model while the
  analytic capsule constraint is active. Model mismatch in capsule radii,
  torso geometry, or exclusions directly affects safety.

  \item \textbf{Effort cost is not applied torque.}
  The virtual PD effort omits actuator-force saturation. Large errors can
  therefore make its numerical penalty differ from MuJoCo's applied force.

  \item \textbf{Several quantities are soft-only.}
  The box target, box uprightness, base tracking, IK reachability during MPPI,
  and EE contact are objectives or references, not formal dynamic constraints.

  \item \textbf{Real-time performance must be measured.}
  Each update evaluates $30\times40$ contact-dynamics steps and rollout
  sensors, while CEM occasionally evaluates 90 IK/collision candidates on a
  background thread. Sustained $100\,\mathrm{Hz}$ operation depends on the
  target CPU and thread contention.
\end{enumerate}

%=============================================================================
\section{Recommended Validation}
%=============================================================================

Useful closed-loop measurements include:
\begin{enumerate}
  \item mean, 95th, and 99th percentile MPPI update time and the 10 ms
        deadline-miss rate, with and without concurrent CEM,
  \item CEM latency, IK residual, accepted/fallback/rejected plan counts, and
        stale-plan discard count,
  \item box $x/y$ error, path length, speed, yaw, tilt, and completion time,
  \item EE-to-moving-target error, engagement time, and measured EE--box
        contact force,
  \item minimum arm--torso clearance, valid-rollout fraction, and all-invalid
        MPPI fallback count,
  \item base tracking error, roll/pitch, foot contacts, and slip,
  \item gait/IK nominal, selected actions, warm-start correction, and actuator
        saturation rate,
  \item one-step error between predicted and executed robot and box states.
\end{enumerate}

Useful ablations include disabling the CEM base planner, removing the
penetration switch, replacing the geometric engagement test with contact
feedback, removing the tilt violation penalty, disabling terminal EE
emphasis, and comparing \code{stance\_hold} noise scale $0.6$ against $1.0$.

%=============================================================================
\section{Pipeline Summary}
%=============================================================================

\begin{equation*}
\boxed{
\begin{array}{c}
\text{Current robot--box state }\hat{\bm{x}}_t\\
\downarrow\\
\text{Box-to-goal direction + moving rear-face EE target}\\
\downarrow\\
\text{Asynchronous kinematic CEM base-pose reference}\\
\downarrow\\
\text{Rate-limited base command + height-conditioned gait + arm IK}\\
\downarrow\\
\text{Whole-body nominal action + shifted MPPI correction}\\
\downarrow\\
\text{Mode-selected cubic-spline sampling }(30\times40\times16)\\
\downarrow\\
\text{Parallel MuJoCo robot--box rollouts + FK sensors}\\
\downarrow\\
\text{Robot/effort + box/tilt + EE/contact + clearance cost}\\
\downarrow\\
\text{Hard arm--torso rejection + exponential weighting}\\
\downarrow\\
\text{Execute the first desired joint-position action}
\end{array}}
\end{equation*}

The defining design choice is the separation between a low-dimensional
kinematic CEM reference planner and a full-dynamics MPPI controller. CEM moves
the base reference toward a posture from which the evolving contact point is
reachable; MPPI then searches whole-body desired joint-position trajectories
whose MuJoCo rollouts actually stabilize the robot and move the free box.

\end{document}
